\section{15.1 모델기반 RL 개관: Dyna-Q 재조명과 ``공짜 모델''}
\begin{frame}{AlphaGo에서 MCTS로}
  \begin{itemize}
    \item 2016년: 딥마인드 \kb{AlphaGo}가 세계 최정상급 프로 이세돌을 꺾음
    \item 바둑판의 국면 수는 \textbf{우주의 원자 수보다 많다}고 알려져
      있어서 ``가치 근사''(DQN, Week 9)만으로는 부족
    \item AlphaGo의 핵심 무기 중 하나 = \kb{몬테카를로 트리 탐색
      (Monte Carlo Tree Search, MCTS)}
    \item Week 7.3에서 맛본 ``모델을 활용한 계획''을, 명확한 규칙이 있는
      게임에서 훨씬 강력하게 확장한 형태
  \end{itemize}
  \vskip0.6em
  \begin{block}{이 장의 위치}
    지금까지의 ML2(DQN · PPO · 모방학습)는 대부분 \textbf{model-free}:
    시행착오로 배움. 이번 장은 그 반대편 --- \textbf{환경의 규칙이
    주어졌을 때 그 모델을 이용해 미리 계획}하고 결정하는 방법.
    이 두 축이 합쳐져야 AlphaGo/AlphaZero가 완성된다.
  \end{block}
\end{frame}
